\section{摘要}\label{sec:摘要}

本報告旨在從不同角度了解AI對教育體系之影響，(一)各國指引：此部分整理現有各國政府教育部們釋出的AI融入教學指引，並進一步進行比較。(二)教育現場的調查：第二部分對台灣教育現場的師生進行初步調查，以問卷蒐集120位國中九年級學生及16位中小學教師對AI融入教學態度與想法，進一步分析與討論。(三)文獻回顧：透過回顧AI融入教學的實徵研究，來了解對學習所帶來的影響有哪些。(四)綜合觀點：在這些結果中我們如何看待AI是否應融入教學，並進一步提出未來建議。

\section{議題簡介：時事案例}\label{sec:時事案例}

二〇二六年四月，美國普渡大學（Purdue University）資訊工程學系爆發了一起震驚全美學術界的事件。該系核心必修課程「CS240 系統程式設計（Systems Programming）」的授課教師 Jeffrey Turkstra 副教授，透過其參與開發並於 2022 年發表於電腦科學教育領域頂尖國際會議 SIGCSE 的程式碼追蹤系統「EnCourse」\cite{rodriguez2022tracking}，向班上超過兩百名學生發出學術誠信違規（Academic Integrity Violation）的指控電子郵件\cite{arun2026ai,lafayette2026cheating}。

這封電子郵件的時機與內容引發了極大的爭議。Turkstra 選擇在該學期的「退選截止日（Withdrawal Deadline）」當天寄出此信，要求被指控的學生填寫一份 Google 表單，自行坦白在哪幾份作業中使用了生成式 AI 工具，並警告若未於限期內回覆，可能面臨課程成績不及格（F）及學務處（Office of the Dean of Students）的進一步調查與懲處\cite{arun2026ai,lafayette2026cheating}。更引發學生恐慌的是，信中並未告知每位學生具體是哪一份作業觸發了系統的警報，學生必須在缺乏具體資訊的情況下自行認罪，形同將舉證責任完全轉嫁給學生\cite{arun2026ai,bailey2026cheating}。

郵件在學生社群中引發了大規模恐慌。根據《Lafayette Journal \& Courier》（USA TODAY Network）的報導，超過半數的學生選擇在截止日前緊急退選課程\cite{lafayette2026cheating}。消息迅速在網路社群擴散，累積了數千則討論，引起了校園報紙《The Purdue Exponent》及主流新聞媒體的關注\cite{arun2026ai,lafayette2026cheating}。在強烈的輿論壓力與學務處的行政介入下，Turkstra 於四月二十日的課堂上宣布了全面的政策撤回：過往所有作業的 AI 作弊指控一律撤銷，認罪表單收集的資料全數作廢，已退選的學生可無條件重新加選\cite{arun2026ai}。

Turkstra 在課堂上展示了課程數據，指出使用 AI 完成作業的學生在實體閉卷考試中的成績，比親自完成作業的學生平均低了百分之十至十五，佐證其在基礎課程中限制 AI 使用的教學立場\cite{arun2026ai}。然而，全面撤銷指控也帶來了新的道德困境。Turkstra 坦言：「我們有許多投入大量心血、沒有作弊的學生，現在卻只能眼睜睜看著部分作弊的同學安然脫身且不用承擔任何後果。對此，我無法給出一個滿意的答案。」\cite{arun2026ai}

這起事件精準地折射出生成式 AI 對教育體系帶來的多重衝擊：在技術層面，自動化偵測工具的準確性與誤判問題亟待解決；在教學層面，基礎課程中應否禁止使用 AI 工具、以及如何重新設計評量方式，成為教育工作者必須面對的核心課題；在制度層面，程序正義與學術誠信之間的平衡，在 AI 時代面臨前所未有的挑戰。這些問題，也正是本報告希望從不同角度深入探討的焦點。

\section{現況分析：政府的指引}\label{sec:政府指引}

面對生成式 AI 帶來的教育衝擊，各國政府與國際組織紛紛制定政策指引，試圖在擁抱科技創新與維護學術品質之間取得平衡。聯合國教科文組織（UNESCO）於二〇二四年正式發布了《學生人工智慧素養架構》（AI Competency Framework for Students），為全球教育體系提供了首份標準化的 AI 能力參考藍圖\cite{unesco2024competency}。根據 UNESCO 的調查，截至二〇二二年，全球僅有十一個國家正式制定並核定 K-12 AI 課程綱要，另有四個國家正在研發中\cite{unesco2022k12}，凸顯了教育體系採用新技術的速度遠落後於科技發展，已成為全球性的教育危機。UNESCO 的素養架構並未將 AI 教育窄化為單純的程式設計或軟體操作教學，而是建構了一套由「四個核心維度」與「三個漸進階層」交織而成的十二項素養矩陣\cite{unesco2024competency}。四個維度涵蓋了「以人為本的心智模式」、「人工智慧倫理」、「AI 技術與應用」以及「AI 系統設計」；三個漸進階層則從「理解（Understand）」出發，經由「應用（Apply）」，最終達到「創造（Create）」的精熟水準。架構的核心精神在於強調人類主體性（Human Agency）與當責性（Accountability）：即使 AI 系統提供了看似權威的建議，人類依然擁有最終的決策權，並必須為 AI 產出的內容承擔完全責任\cite{unesco2024competency}。

在國內方面，臺灣教育部因應生成式 AI 的快速發展，將《中小學使用生成式人工智慧注意事項》更新至 2.1 版，並依據不同學習階段的認知發展特徵，編製了兩套學習應用手冊：針對國小學生的《我和 AI 一起學》以及針對國高中學生的《駕馭 AI，洞察未來：數位公民的必修課》\cite{govtw2024aiguide}。國小版手冊以「5W1H」架構引導學童建立對 AI 的基礎認知與自我保護意識，強調不洩露個人隱私資料、辨識 AI 幻覺（Hallucination）等核心觀念；國高中版則將教育目標提升至培養「AI 共學力」，強調操作技能精進、批判性內容查核與倫理反思\cite{govtw2024aiguide}。在高等教育層面，各大專校院也紛紛發布校內指引。以國立臺灣師範大學為例，其教學發展中心公布的《生成式 AI 之學習應用及參考指引》明確列出五大核心倫理面向（公平性、隱私與資料保護、透明與可解釋性、責任歸屬以及人機界線與原創性），並要求師生在使用 AI 時必須進行自我揭露，清楚標註 AI 參與的範疇，同時鼓勵各教學單位依學門特性自訂更細緻的規範\cite{ntnu2025ai}。

跨國比較則揭示了各國在政策取向上的顯著差異。日本文部科學省於二〇二四年十二月發布了《初等中等教育階段生成式 AI 使用指引（Ver. 2.0）》（初等中等教育段階における生成AIの利活用に関するガイドライン），展現了最為審慎的立場\cite{mext2024genai}。該指引強調不應統一禁止或強制使用 AI，而是應評估其使用是否有助於課程綱要所揭示的素養與能力培育，但同時也明確警告，過度依賴 AI 可能導致學生對 AI 生成答案的「盲目接受」，進而對「思考力、判斷力與表達力」的發展造成不利影響\cite{mext2024genai}。

相較之下，英國教育部（DfE）的 \textit{Generative Artificial Intelligence (AI) in Education} 政策文件則採取了「雙軌制」策略：一方面大力鼓勵教師端使用 AI 以減輕行政負擔（如課程規劃、回饋生成），甚至資助 Oak National Academy 開發了 AI 教學助理工具「Aila」\cite{dfe2023genai}，另一方面則要求學校在使用任何 AI 工具前必須進行全面的風險評估，嚴格遵守資料保護法與《在學校中維護兒童安全》（\textit{Keeping Children Safe in Education}）的法定責任，並關注深度偽造（Deepfakes）等新興風險\cite{dfe2023genai}。

至於美國，教育部教育技術辦公室於二〇二三年五月發布的 \textit{Artificial Intelligence and the Future of Teaching and Learning} 報告則將視角拉高至社會公平層面\cite{usedu2023ai}。該報告明確提出 AI 應「增強人類智能而非取代人類」，強調必須保持「人類在迴路（Human-in-the-loop）」的核心原則，並聚焦管理「演算法偏見」的風險，因為 AI 系統的開發過程可能導致模式偵測的偏差與自動化決策的不公，進而邊緣化弱勢群體\cite{usedu2023ai}。

新加坡則展現了最為積極的擁抱姿態。根據經濟合作暨發展組織（OECD）二〇二四年《教與學國際調查》（TALIS）的數據，新加坡高達百分之七十五的中學教師已在教學中常態性使用 AI，遠超 OECD 百分之三十六的平均水準\cite{moe2025singapore}。新加坡教育部提出了「高信任、高接觸、高科技」的教育架構，將 AI 定位為實現「有效班級規模為一」（即高度個人化學習）的關鍵工具，同時堅持人類教師在社會情感教育中的不可替代性\cite{moe2025singapore}。在高教端，新加坡國立大學（NUS）制定了嚴謹的生成式 AI 使用規範，確立了「使用者須對 AI 生成的所有內容負完全責任，而非 AI 工具本身」的核心原則，並要求學生在提交作業時必須揭露使用的 AI 工具名稱、使用方式及具體貢獻範圍\cite{nus2024integrity}。

綜觀各國政策，可以發現若干高度一致的共識：第一，AI 不應取代人類的判斷與決策，教育者與學習者必須保持「人類在迴路」的主體地位；第二，所有允許 AI 參與的場景，皆以「強制揭露」與「人類負完全責任」為前提；第三，各國普遍認知到，教育體系不應僅停留在「防堵 AI」的消極立場，而應積極建構系統性的 AI 素養教育。然而，在開放程度與監管強度上，各國仍有顯著分歧，這也反映了不同文化底蘊、教育體制與國家科技發展戰略的深層差異。這些官方指引為教育現場提供了原則性架構，但實際的執行成效與師生的真實態度，仍需要透過實證調查來進一步驗證。

\section{實際調查：問卷設計與發放}\label{sec:問卷設計}

本報告旨在探討教學現場師生對生成式人工智慧（AI）融入教育的看法，因此採用問卷調查法進行資料蒐集。根據生成式AI於教育領域的應用現況及相關文獻資料，自行設計教師版與學生版兩份問卷，以了解師生對AI工具的使用情形、態度及影響認知。

本研究以\textbf{120位國中九年級學生}及\textbf{16位中小學教師}為調查對象。學生問卷主要針對學生在學習過程中使用生成式AI工具的情形進行調查；教師問卷則聚焦於教師在教學現場使用AI工具的經驗，以及對AI融入教育的看法。問卷採匿名方式填寫，以提高受試者作答意願及資料真實性。

學生問卷內容包含五個面向：（一）AI工具使用頻率與情境、（二）AI使用方式、（三）AI依賴程度、（四）學習行為是否改變、（五）AI倫理與使用態度，共24題。透過這些題項了解學生使用AI工具的實際情況及其對學習產生的影響。

教師問卷則包含背景資料調查、AI使用經驗、學生AI使用觀察、教師對AI融入教育的態度，以及開放式意見回饋等部分。內容涵蓋教師使用AI進行備課、教材設計、行政工作等情形，並探討教師對AI對學生學習、教學方式及未來教育發展之看法。

問卷回收後，研究團隊將資料進行整理與統計分析，藉以了解師生對AI融入教學的認知、使用現況及態度差異，作為未來教育現場推動AI應用之參考依據。

\section{分析結果}\label{sec:分析結果}

\subsection{學生的態度}\label{sec:學生態度}

\subsubsection{1. 構面：使用方式}

\begin{longtable}{l c c}
\toprule
題目 & 平均數 & 標準差 \\
\midrule
\endfirsthead
\toprule
題目 & 平均數 & 標準差 \\
\midrule
\endhead
\midrule
\multicolumn{3}{r}{續下頁} \\
\endfoot
\bottomrule
\endlastfoot
5. 我會先自己思考，再詢問AI & 4.00 & 1.000 \\
6. 我會直接使用 AI 提供的答案 & 2.68 & 1.033 \\
7. 我會修改 AI 產生的內容 & 3.66 & 1.203 \\
8. 我會確認 AI 回答是否正確 & 3.75 & 1.147 \\
9. 我會請 AI 解釋不懂的概念 & 4.09 & 0.974 \\
10. 我會透過反覆提問讓 AI 提供更好的答案 & 3.72 & 1.221 \\
\end{longtable}

在使用方式構面中，可以從結果分析看出「先自己思考再詢問 AI」（Q5，M = 4.00）以及「請 AI 解釋不懂的概念」（Q9，M = 4.09）兩題平均數較高，顯示\textbf{學生較常將 AI 用於概念理解與學習輔助}。

\subsubsection{2. 構面：依賴程度}

\begin{longtable}{l c c}
\toprule
題目 & 平均數 & 標準差 \\
\midrule
\endfirsthead
\toprule
題目 & 平均數 & 標準差 \\
\midrule
\endhead
\midrule
\multicolumn{3}{r}{續下頁} \\
\endfoot
\bottomrule
\endlastfoot
11. 沒有 AI 時，我會不知道如何開始作業 & 1.61 & 0.922 \\
12. 我習慣先詢問 AI，而不是自己查資料 & 2.22 & 1.206 \\
13. 我認為 AI 比自己搜尋資料更方便 & 3.17 & 1.174 \\
14. 使用 AI 後，我較少主動思考問題 & 2.13 & 1.094 \\
15. 我越來越依賴 AI 協助學習 & 2.18 & 1.110 \\
\end{longtable}

在依賴程度構面中，各題的平均數普遍偏低，尤其是 Q11 的分數接近極小值。反映出多數學生並未對 AI 產生過度依賴，仍自認具備獨立規劃並完成學習任務的核心能力。

\subsubsection{3. 構面：學習行為改變}

\begin{longtable}{l c c}
\toprule
題目 & 平均數 & 標準差 \\
\midrule
\endfirsthead
\toprule
題目 & 平均數 & 標準差 \\
\midrule
\endhead
\midrule
\multicolumn{3}{r}{續下頁} \\
\endfoot
\bottomrule
\endlastfoot
16. AI 改變了我的學習方式 & 2.65 & 1.222 \\
17. 使用 AI 後，我更有效率地完成學習任務 & 3.46 & 1.023 \\
18. 使用 AI 後，我花較少時間查資料 & 3.24 & 1.133 \\
19. 使用 AI 後，我更願意自主學習 & 2.92 & 1.005 \\
20. AI 讓我更容易理解困難內容 & 3.29 & 1.075 \\
\end{longtable}

在此構面中，Q17 與 Q20 的平均數落在 3.2 至 3.5 之間，屬於中立區間。這顯示學生對於「AI 是否顯著改變其學習效率或理解困難內容」的看法尚不一致，整體態度偏向中立，並未形成壓倒性的普遍認同。

\subsubsection{4. 構面：倫理與使用態度}

\begin{longtable}{l c c}
\toprule
題目 & 平均數 & 標準差 \\
\midrule
\endfirsthead
\toprule
題目 & 平均數 & 標準差 \\
\midrule
\endhead
\midrule
\multicolumn{3}{r}{續下頁} \\
\endfoot
\bottomrule
\endlastfoot
21. 我認為直接繳交 AI 生成內容是不恰當的 & 3.67 & 1.288 \\
22. 我會擔心 AI 提供錯誤資訊 & 4.10 & 0.982 \\
23. 我認為學生應學習如何正確使用 AI & 4.18 & 0.908 \\
24. 我認為 AI 對學習的幫助大於壞處 & 3.45 & 0.963 \\
\end{longtable}

在倫理與使用態度構面中，可以從題目「學生普遍認為應學習如何正確使用 AI」（Q23，M = 4.18），並對「AI 可能產生錯誤資訊抱持一定程度的擔憂」（Q22，M = 4.10），表示學生普遍認為\textbf{AI 仍應作為輔助工具，而不是直接取代自己的學習成果}。

\subsubsection{結論}

整體而言，可以從表單結果看出學生普遍認同 AI 能協助學習與提升效率，尤其在理解困難內容、完成學習任務以及查詢資料方面。

此外，多數學生表示在使用 AI 前仍會先自行思考，並會確認或調整 AI 所提供的內容，顯示學生傾向將 AI 作為學習輔助工具，而非直接依賴 AI 完成所有學習任務。

在倫理與使用態度方面，學生普遍認同應學習如何正確使用 AI，並對 AI 可能產生錯誤資訊抱持一定程度的警覺，也認為直接繳交 AI 生成內容並不恰當，顯示學生對 AI 使用仍具備一定的學術倫理道德。

整體而言，本分析結果可以看出，生成式 AI 已逐漸成為學生學習過程中的輔助工具。學生大多肯定 AI 在學習上的便利性與幫助，但同時也對其使用方式與資訊正確性保持一定程度的注意與判斷。
